\documentclass[10pt,a4paper]{article}
\usepackage[utf8]{inputenc}
\usepackage[german]{babel}
\usepackage[margin=1.5cm]{geometry}
\usepackage{titlesec}
\usepackage{tcolorbox}
\usepackage{multicol}
\usepackage{booktabs}
\usepackage{csquotes}

% Design-Anpassungen für Fantasy-Stile
\titleformat{\section}{\large\bfseries\sffamily\uppercase}{}{0em}{}[\titlerule]
\titlespacing{\section}{0pt}{12pt}{6pt}

\tcolorboxenvironment{minipage}{
    colback=gray!5,
    colframe=gray!40,
    arc=2mm,
    boxrule=0.5mm,
    boxsep=2mm
}

\begin{document}

% Basisdaten
\begin{center}
    {\Huge\bfseries\sffamily Valerian \enquote{Silvertongue} Vance} \\[0.2cm]
    {\large Ein Pathfinder Charakter-Almanach}
\end{center}

\vspace{0.3cm}

\noindent
\begin{tabular*}{\textwidth}{@{\extracolsep{\fill}}llll}
    \textbf{Volk/Abstammung:} Mensch (Chelaxier) & \textbf{Klasse:} Schurke & \textbf{Stufe:} 1 \\
    \textbf{Gesinnung:} Chaotisch Gut & \textbf{Gottheit/Glaube:} Cayden Cailean & \textbf{Heimat:} Kintargo \\
\end{tabular*}

\vfill
\hrule
\vfill

\begin{multicols}{2}

\section{Das Äußere \& Erster Eindruck}
\begin{itemize}
    \item \textbf{Erscheinung:} Vernarbte Hände, wettergegerbtes Gesicht, trägt ein markantes, abgewetztes Lederwams.
    \item \textbf{Verhalten:} raue Stimme, anspruchsvoller Weggefährte, stets wachsam.
    \item \textbf{Besonderes Merkmal:} TODO
\end{itemize}

\section{Werdegang \& Hintergrund}
Valerian wurde im Herzen von Kintargo geboren, genauer gesagt am Rande des berüchtigten Viertels Devil’s Nursery (der Teufelsgärtnerei) und den nahegelegenen Docks am Yolubilis-River. Die Mischung aus den bitterarmen Elendsvierteln der Tieflinge und dem geschäftigen, schmutzigen Treiben des Hafens war seine Schule. Er wuchs ohne den Prunk der Adelsfamilien auf, lernte aber schnell, dass man in Kintargo entweder Raubtier oder Beute ist.

Da es in Kintargo seit Längerem keine zentralisierte, organisierte Diebesgilde mehr gibt, musste Valerian als freischaffender „Spezialist“ überleben. Er schloss sich weder den Red Jills noch den River Talons an, da diese Banden keine neuen Mitglieder aufnahmen. Stattdessen perfektionierte er das Handwerk des klassischen Einbrechers: Er knackte die Tresore korrupter Händler und stahl die Luxusgüter der Thrune-treuen Oberschicht.
Sein Markenzeichen war jedoch nicht nur seine Fingerfertigkeit, sondern sein Mundwerk: Mehr als einmal entkam er den Stadtwachen (Dottari), indem er sich mit Charme, erfundenen Identitäten oder Bestechung aus brenzligen Situationen herausredete.

Valerian war nicht immer ein verbitterter Krimineller. Er hatte eine Schwester, Elysia, eine talentierte Schneiderin, die im noblen Teil der Stadt arbeitete. Vor zwei Jahren wurde Elysia jedoch Opfer der drakonischen Gesetze des Hauses Thrune. Sie wurde fälschlicherweise beschuldigt, verbotene Flugblätter der Silver Ravens (einer alten Rebellengruppe) besessen zu haben. Ohne faires Verfahren wurde sie von den Inquisitoren verschleppt. Valerian versuchte verzweifelt, in das Gefängnis einzubrechen, um sie zu befreien – sein erster großer Rückschlag. Das Schloss war magisch gesichert, er scheiterte, entkam nur knapp und Elysia wurde kurz darauf exekutiert. Dieser Verlust brach ihm das Herz und machte aus dem einfachen Dieb einen Mann mit tiefem Hass auf das Regime.

\section{Charakterzüge \& Philosophie}
\begin{itemize}
    \item \textbf{Motivation:} TODO
    \item \textbf{Ziele:} TODO
\end{itemize}

\section{Schwächen}
\begin{itemize}
    \item \textbf{Grenzenlose Arroganz (Überheblichkeit:)} Valerian weiß, wie gut er ist - und das ist sein größtes Problem. Er ist fest davon überzeugt, dass er der klügste im Raum ist. Diese Arroganz führt dazu, dass er Warnungen von Gefährten oft abtut. Er glaubt sich aus jeder Situation herausreden oder herauswinden zu können.
    \item \textbf{Rachsucht:} Der Verlust seiner Schwester Elysia hat eine tiefe, schwärende Wunde hinterlassen. Sobald es um den \itextit{Order of the Rack} geht verliert der Schurke schnell seine Fassung und wird unvorsichtig.
    \item \textbf{Angst vor echter Nähe:} Valerian benutzt seinen Charme als Schutzschild. Er flirtet scherzt und verbrüdert sich schnell, aber das ist alles nur Fassade. Er lässt niemanden emotional an sich heran.
\end{itemize}

\section{Vorlieben}

\begin{itemize}
    \item \textbf{Luxus (als Diebesgut):} Er stiehlt nicht nur zum Überleben, sondern auch aus Genuss. Er liebt die feinen Dinge, die das Haus Thrune den normalen Bürgern verwehrt. Er hat eine Schwäche für teuren cheliaxischen Wein (aus herzoglichen Kellern), feine Seidentücher und maßgeschneiderte Kleidung.
    \item \textbf{Kunst des Schlossknackens:} Er betrachtet Schlösser wie ein Rätsel, die der Erfinder speziell für ihn hinterlassen hat. Je komplexer, desto größer ist seine Freude, wenn das Klicken ertönt.
    \item \textbf{Perfekter Abgang:} Weil er so arrogant ist reicht es ihm oft nicht, einfach unbemerkt zu entkommen. Er verspürt den Drang, eine Visitenkarte zu hinterlassen oder dem Opfer eine spöttische Nachricht zu schreiben. Dieses Geltungsbedürfnis gefährdet nicht nur ihn selbst, sondern im Verlauf der Kampagne auch die Geheimhaltung der gesamten Rebellion.
\end{itemize}


\section{Glaube}
Nach dem Tod seiner Schwester verlor Valerian den Glauben an eine gerechte Ordnung (wie sie Abadar predigt) komplett. Stattdessen fand er Trost im Verborgenen. Er betet heimlich zu Cayden Cailean, dem Gott der Freiheit und des Aufbegehrens. Da die offizielle Kirche des Gottes im Bürgerkrieg niedergebrannt und durch den düsteren Tempel von Zon-Kuthon (Shadowsquare) ersetzt wurde , sucht Valerian die geheimen Schreine des „Trunkenen Gottes“ auf, die in den Hinterzimmern von Kintargos Tavernen und Bordellen versteckt sind. Er verbirgt seine religiösen Symbole geschickt vor den Augen der Obrigkeit.


\section{Geheimnisse}
Bei einem seiner Einbrüche in ein Regierungsgebäude (House of Truth and Clarity) kurz vor der Verhängung des Kriegsrechts stahl Valerian kein Gold, sondern Dokumente. Er besitzt eine geheime, unzensierte Akte aus der Zeit vor der Thrune-Zensur. Diese Dokumente beweisen, dass die Historiker von Haus Thrune die Geschichte der Stadt massiv gefälscht haben. Er versteckt diese Papiere an einem sicheren Ort, da ihr Besitz sein sicheres Todesurteil durch Barzillai Thrune bedeuten würde.

\section{Wichtige Beziehungen}
\begin{itemize}
    \item \textbf{Elysia (Tod):} Seine geliebte Schwester. Wurde von den Inquisitoren hingerichtet.
    \item \textbf{Geoff Tanessen} Hat das Todesurteil von Elysia unterzeichnet. Valerion hat sich geschworen Rache zu nehmen, komme was wolle.
\end{itemize}

\end{multicols}

% Sonstige Tipps zur Spielweise
\begin{tcolorbox}[title=Tipps vor der Kampagne, colframe=gray!60, colback=white, arc=1mm]
    \small
    Wenn ein Plan der Gruppe fehlschlägt, weil Valerian ein Schloss nicht schnell genug knackt, würde er niemals zugeben, dass es sein Fehler war. Er würde eher behaupten: „Das Schloss war manipuliert, eine billige imperiale Fehlkonstruktion! Kein Meisterdieb hätte das unter fünf Minuten geschafft.“ Wenn er jemanden überzeugen muss (Diplomatie-Wurf), setze auf eine herablassende, aber unwiderstehliche Eleganz – er gibt den Leuten das Gefühl, dass es ein Privileg ist, mit ihm zu sprechen.
\end{tcolorbox}

\end{document}
